\section{Introduction}

Artificial intelligence is now embedded in production planning, service delivery, logistics, and knowledge work. The policy debate has therefore shifted from whether AI can improve specific tasks to whether it can raise aggregate productivity and growth in a sustained way \cite{bresnahanTrajtenberg1995,cominMestieri2018}. Although firm-level and sectoral studies often document efficiency gains, cross-country macro evidence remains incomplete and methodologically fragmented \cite{brynjolfssonRockSyverson2019,noyZhang2023,brynjolfssonLiRaymond2023}. Three empirical frictions are central: AI adoption is measured with heterogeneous proxies, cross-country data quality differs materially, and macro outcomes are jointly determined by multiple evolving structural factors \cite{imf2024,oecd2023ai,goldfarbTucker2019}.

These constraints matter because they create a gap between optimistic policy narratives and the strength of available evidence. In particular, large aggregate claims can exceed what observational macro panels can credibly identify without strong assumptions. A central objective of this paper is therefore to narrow the gap between empirical ambition and inferential discipline.

This paper addresses two questions: whether higher AI adoption is associated with higher TFP within countries over time, and whether a similar pattern appears in near-term GDP growth. We estimate panel models centered on within-country variation and evaluate stability under pre-specified robustness and sensitivity checks.

The focus on within-country variation is deliberate. Country fixed effects absorb persistent structural differences (for example, long-run institutions, geography, and industrial mix), while supplementary specifications evaluate the extent to which inference is sensitive to alternative assumptions about common shocks, outliers, timing, and dependence structures.

Our contribution is threefold. We provide updated cross-country evidence on AI and productivity using a harmonized panel and transparent variable construction. We report a structured robustness portfolio rather than relying on a single preferred estimator, allowing stability to be assessed across plausible alternatives. We also implement a reproducible pipeline that links data validation, estimation, tables, and manuscript text, reducing reporting inconsistency risk.

An additional contribution is translational. The paper connects micro-level AI productivity narratives to macro-level policy interpretation without overstating identification. Rather than treating strong task-level gains as direct evidence of aggregate transformation, it tests whether comparable country-level associations are visible, stable, and economically plausible in available panel data.

This framing matters for applied interpretation. Development strategies are often formulated under uncertainty about effect timing and scale. By centering within-country variation, reporting model-dependent sample limits, and applying pre-specified robustness checks, the manuscript offers a usable evidence baseline while keeping causal claims proportionate to design strength.

Previewing results, we find a positive and statistically precise AI-TFP association in the baseline model, with directional stability across two-way fixed effects, trimming, lag specifications, and alternative inference choices. The GDP growth association is positive but notably smaller, consistent with a staged transmission process in which productivity responses precede broader macro growth responses.

We emphasize from the outset that these results should be interpreted as conditional associations. The design does not by itself establish final causal parameters, especially under potential reverse causality and time-varying omitted factors.

The remainder of the paper proceeds as follows. Section 2 situates the study in related literature. Section 3 describes data and variable construction. Section 4 presents the empirical strategy. Section 5 reports results. Section 6 discusses interpretation and limitations. Section 7 concludes.
